\section{Limitações}
\label{sec:limitacoes}

\textbf{Cobertura.} O ano de 2026 é parcial e os índices podem mudar. Não houve pré-registro público do protocolo. Também não foram realizadas busca de citações para frente ou para trás nem busca estruturada de literatura cinzenta, e a cobertura permaneceu condicionada às fontes e aos metadados indexados, apesar da ausência de filtros de idioma ou tipo documental.

\textbf{Seleção e codificação.} As decisões de triagem e codificação foram conduzidas por um único avaliador, sem segundo revisor independente e sem medida de concordância interavaliadores, e o codebook e termos amplos, como \textit{framework}, incorporam escolhas analíticas. Presença lexical, variável mensurada, critério aplicado e resultado empírico não são equivalentes; ODS, SDG e ESG podem estar ausentes mesmo quando seus componentes aparecem.

\textbf{Extração e força de evidência.} O núcleo de 104 não foi submetido integralmente a elegibilidade e extração em texto completo. Trinta estudos com PDF disponível foram posteriormente lidos em tarefas pontuais e 14 forneceram evidências específicas. A unidade de análise quantitativa é o registro bibliográfico consolidado. Ela não garante independência entre publicações relacionadas nem representa força de evidência acumulada; a classificação de desenho deixou 23,1\% sem sinal suficiente e nenhum instrumento de risco de viés foi aplicado, o que impede comparar qualidade metodológica ou recomendar superioridade.

\textbf{Bibliometria.} Pela incompletude de autores e pela ausência de afiliações e referências citadas nos 372 registros (Seção~\ref{sec:bibliometriaampliada}), colaboração e acoplamento bibliográfico não foram executados; mapa temático, evolução e coocorrências permanecem explorações documentais, sem direção causal ou peso científico.

\textbf{Protocolo.} Indicadores e regras formam especificação operacional candidata: pesos, limiares, função de agregação e sensibilidade ainda não foram validados nem executados com dados reais, permanecendo dependentes de validação de conteúdo, participação institucional e teste de robustez.
